%%%%%%%%%%%%%%%%%%%%%%%%%%%%%%%%%%%%%%%%%
% baposter Landscape Poster
% LaTeX Template
% Version 1.0 (11/06/13)
%
% baposter Class Created by:
% Brian Amberg (baposter@brian-amberg.de)
%
% License:
% CC BY-NC-SA 3.0 (http://creativecommons.org/licenses/by-nc-sa/3.0/)
%
%%%%%%%%%%%%%%%%%%%%%%%%%%%%%%%%%%%%%%%%%

%----------------------------------------------------------------------------------------
% PACKAGES AND OTHER DOCUMENT CONFIGURATIONS
%----------------------------------------------------------------------------------------

\documentclass[landscape,a0paper,fontscale=0.285]{baposter} % Adjust font scale/size

\usepackage{graphicx} % Required for including images
\graphicspath{{figs/}} % Directory in which figures are stored

\usepackage{amsmath} % For typesetting math
\usepackage{amssymb} % Adds new symbols to be used in math mode
\usepackage{mathtools}
\usepackage{url}
\usepackage[numbers]{natbib}

\usepackage{booktabs} % Top and bottom rules for tables
\usepackage{enumitem} % Used to reduce itemize/enumerate spacing
\usepackage{palatino} % Use the Palatino font
\usepackage[font=small,labelfont=bf]{caption} % Required for specifying captions to tables and figures

\usepackage{multicol} % Required for multiple columns
\setlength{\columnsep}{1.5em} % Slightly increase the space between columns
\setlength{\columnseprule}{0mm} % No horizontal rule between columns

\newcommand{\compresslist}{ % Define a command to reduce spacing within itemize enumerate environments, this is used right after \begin{itemize} or \begin{enumerate}
\setlength{\itemsep}{1pt}
\setlength{\parskip}{0pt}
\setlength{\parsep}{0pt}
}

\definecolor{lightblue}{rgb}{0.145,0.6666,1} % Defines color of content box headers

\begin{document}

\begin{poster} {
headerborder=closed, % Adds a border around the header of content boxes
colspacing=1em, % Column spacing
bgColorOne=white, % Background color for the gradient on the left side of the poster
bgColorTwo=white, % Background color for the gradient on the right side of the poster
borderColor=lightblue, % Border color
headerColorOne=black, % Background color for the header in the content boxes (left side)
headerColorTwo=lightblue, % Background color for the header in the content boxes (right side)
headerFontColor=white, % Text color for the header text in the content boxes
boxColorOne=white, % Background color of the content boxes
textborder=roundedleft, % Format of the border around content boxes, can be: none, bars, coils, triangles, rectangle, rounded, roundedsmall, roundedright or faded
eyecatcher=true, % Set to false for ignoring the left logo in the title and move the title left
headerheight=0.1\textheight, % Height of the header
headershape=roundedright, % Specify the rounded corner in the content box headers, can be: rectangle, small-rounded, roundedright, roundedleft or rounded
headerfont=\Large\bf\textsc, % Large, bold and sans serif font in the headers of content boxes
%textfont={\setlength{\parindent}{1.5em}}, % Uncomment for paragraph indentation
linewidth=2pt % Width of the border lines around content boxes
}
%----------------------------------------------------------------------------------------
% TITLE SECTION
%----------------------------------------------------------------------------------------
%
{\includegraphics[height=6.5em]{logo_berkeley.jpg}} % First university/lab logo on the left
{\bf\textit{\LARGE Nonparametric-efficient Causal Mediation Analysis for
    Stochastic Interventions}\vspace{0.01em}} % Poster title
{\textbf{Nima Hejazi, Mark van der Laan, and Iv{\'a}n D{\'\i}az} \\
  \textit{Graduate Group in Biostatistics \& Dept.~of Statistics, UC Berkeley}\\
  \textit{Division of Biostatistics, Dept.~of Healthcare Policy \& Research,
          Weill Cornell Medicine}}
    % Author names and institution
%{\includegraphics[height=6.5em]{logo_weill.jpg}} % Second university/lab logo on the right

%-------------------------------------------------------------------------------
% OVERVIEW
%-------------------------------------------------------------------------------

\headerbox{Overview \& Motivations}{name=overview,column=0,row=0}{

\begin{itemize}\compresslist
  \setlength\itemsep{0.75em}
  \item We present a decomposition of the \textit{population intervention
      effect}, defined through stochastic interventions.
  \item Population intervention effects provide a generalized framework in which
    a variety of interesting causal contrasts can be defined, including effects
    for continuous and categorical exposures.
  \item We propose estimators of these direct and indirect effects:
    \begin{itemize}
      \item substitution (G-computation), re-weighted (IPW), and
      \item efficient estimators based on flexible regression techniques
        (one-step, TMLE).
    \end{itemize}
\end{itemize}
\vspace{0.65cm} % When there are two boxes, some whitespace may need to be added
                % if the one on the right has more content
}

%-------------------------------------------------------------------------------
% INTRODUCTION
%-------------------------------------------------------------------------------

\headerbox{Software Implementation}
{name=introduction,column=1,row=0,bottomaligned=overview}{

\begin{itemize}\compresslist
  \setlength\itemsep{0.5em}
  \item Our efficient estimators are asymptotically linear under a condition
    requiring $n^{1/4}$-consistency of certain regression functions.
  \item The efficient estimators are asymptotically normal with estimable
    variance, thereby allowing for the construction of confidence intervals and
    hypothesis tests.
  \item The \underline{\texttt{medshift} \texttt{R} package}
    \cite{hejazi2019medshift} implements these estimators and leverages
    state-of-the-art machine learning in the procedure.
  \item \texttt{tlverse} plug: \url{https://tlverse.org}
  \item ...
\end{itemize}
}

%-------------------------------------------------------------------------------
% Methodology cont.
%-------------------------------------------------------------------------------

\headerbox{Construction of Nonparametric-Efficient Estimators}{name=results,
  column=2,span=2,row=0}{

\vspace{-0.35em}
\begin{itemize}
\itemsep0pt
\item To avoid entropy conditions on the initial estimators, we rely on
    cross-fitting
    \citep{zheng2011cross, chernozhukov2016double} --- denote by
    $\hat \eta_{j}$ the estimator of $\eta = (g, m, e, \phi)$, obtained by
    training the corresponding prediction algorithm using only data in the
    sample ${\cal T}_j$. Further, let $j(i)$ denote the index of the
    validation set which contains observation $i$. The estimator is thus
    defined as:
    \begin{equation}\label{eq:aipw}
      \hat{\theta}{\delta} = \frac{1}{n} \sum_{i = 1}^n D_{\hat\eta_{j(i)},
          \delta}(O_i) =
        \frac{1}{n}\sum_{i = 1}^n \left\{D^Y_{\hat\eta_{j(i)}, \delta}(O_i) +
        D^A_{\hat\eta_{j(i)}, \delta}(O_i) + D^{Z,W}_{\hat\eta_{j(i)},
        \delta}(O_i) \right\}.
    \end{equation}
\item TMLE!
\end{itemize}
}

%-------------------------------------------------------------------------------
% REFERENCES
%-------------------------------------------------------------------------------

\headerbox{\small References}{name=references,column=2,above=bottom}{
\renewcommand{\section}[2]{\vskip 0.05em} % remove "References" section title
\tiny{ % Reduce the font size in this block
\setlength{\bibsep}{0.25pt}
\bibliographystyle{IEEEtranS}
%\nocite{*} % Insert publications even if they are not cited in the poster
\bibliography{2019_acic}
\compresslist
\vspace{-0.7em}
}
}

%-------------------------------------------------------------------------------
% CONTACT
%-------------------------------------------------------------------------------

\headerbox{\small But wait, there's more!}{name=ack,column=3,aligned=references,above=bottom}{
% This block is as tall as the references block
\begin{itemize}
  \itemsep0.25pt
  \item \textbf{N.~Hejazi}: \texttt{nhejazi@berkeley.edu}; \textbf{M.~van der
      Laan}: \texttt{laan@berkeley.edu}; \textbf{I.~Diaz}:
    \texttt{ild2005@med.cornell.edu}
  \item Pre-print of our original paper: \url{https://arxiv.org/abs/1901.02776}
  \item Iv{\'a}n's talk tomorrow morning!
\end{itemize}
}

%-------------------------------------------------------------------------------
% CONCLUSION
%-------------------------------------------------------------------------------

\headerbox{Results \& Discussion}
{name=conclusion,column=2,span=2,row=0,below=results,above=references}{

\vspace{0.25em}

\begin{multicols}{2}

%\begin{center}
%\vspace*{-0.5cm}
%\includegraphics[scale=0.37]{bias_composed_plot}
%%\captionof{figure}{This figure demonstrates...}
%\end{center}

%\begin{center}
%%\vspace*{-0.49cm}  % CHANGE THIS TO ALIGN IMAGES
%\includegraphics[scale=0.39]{rel_eff_composed_plot}
%%\captionof{figure}{This figure demonstrates...}
%\end{center}

%\vspace{-15pt}

\begin{itemize}
  \itemsep0pt
  \item All estimators approx.~unbiased in large samples; however, inefficient
    TMLE with HAL has bias not converging at $n^{-\frac{1}{2}}$.
  \item Fitting $\Pi$ with HAL or GLM, efficient TMLE has lower variance than
    the inefficient.\\
\end{itemize}

\end{multicols}
}

%-------------------------------------------------------------------------------
% METHODS
%-------------------------------------------------------------------------------

\headerbox{Stochastic Population Intervention (In)Direct Effects}{name=method,
  column=0,span=2,below=overview, bottomaligned=references}{
% This block's bottom aligns with the bottom of the conclusion block

\begin{itemize}
  \itemsep0.75pt
  \item Consider $O = (W, A, Z, Y) \sim P_0 \in \mathcal{M}$, where $W$ is a set
    of baseline covariates, $A$ an intervention, $Z$ a mediator between $A$ and
    outcome, and $Y$ the outcome, with no assumptions on model $\mathcal{M}$.
  \item We may decompose the PIE in terms of a \textit{population intervention
    direct effect (PIDE)} and a \textit{population intervention indirect effect
    (PIIE)}:
    \begin{equation}\label{eqn:pie_decomp}
      \psi(\delta) = \overbrace{\mathbb{E}\{Y(g, q) -
                                      Y(g_\delta, q)\}}^{\text{PIDE}} +
                                    \overbrace{\mathbb{E}\{Y(g_\delta, q) -
                                      Y(g_\delta, q_\delta)\}}^{\text{PIIE}}.
    \end{equation}
  \item The causal parameter is identified by the observed data parameter
    \cite{diaz2019causal}:
    \begin{equation}
      \theta(\delta) = \int m(a, z,
      w)g_\delta(a\mid w)p(z,w) d\nu(a,z,w).
    \end{equation}


  \item Let $\eta = (g, m, e, \phi)$. The efficient influence function for
      $\theta(\delta)$ in the nonparametric model $M$ is
      $D^Y_{\eta, \delta}(o) + D^A_{\eta, \delta}(o) +
      D^{Z, W}_{\eta, \delta}(o) - \theta(\delta)$, where
      \begin{align*}
        D^Y_{\eta, \delta}(o) &= \frac{g_{\delta}(a \mid w)}{e(a \mid z, w)}
                                \{y - m(z,a,w) \};
        D^{Z, W}_{\eta, \delta}(o) = \int m(z, a, w) g_{\delta}(a \mid w)
                                     d\kappa(a),\\
        D^A_{\eta, \delta}(o) &= \frac{g_\delta(a \mid
          w)}{g(a \mid w)}\left\{\phi(a, w) - \int
          \phi(a, w)g_\delta(a\mid w)d\kappa(a)\right\}.
      \end{align*}
      and $\phi_{0}(a, w) = \mathbb{E}\left\{\frac{g(A \mid W)}{e(A \mid Z,
      W)} m(Z, A, W) \mid A = a, W = w \right\}$.
\end{itemize}
}

%-------------------------------------------------------------------------------

\end{poster}
\end{document}
